
% Thuật toán, pseudocode và flowchart	
\chapter{Thuật toán, mã giả và luồng hoạt động}
\subsection{Cơ sở lý thuyết thuật toán AES}

\subsubsection{Mã hóa đối xứng}

Mã hóa khóa đối xứng là một loại phương pháp mã hóa sử dụng một khóa giống nhau để mã hóa, cũng như giải mã dữ liệu.

Mã hóa khối đối xứng bắt đầu với thủ tục sinh khóa, tạo ra một khóa bí mật. Trong quá trình mã hóa, hệ thống dùng các thuật toán 


\subsection{Thuật toán AES}

Thuật toán AES 

\subsubsection{}
Các phép biến đổi khác nhau với kết quả trung gian được gọi là \textbf{Trạng thái}. Trạng thái có thể được biểu diễn dưới dạng byte hình chữ nhật. Một mảng thường bao gồm 4 hàng và số cột, kí hiệu là $N_b$, bằng chiều dài của khối chia cho 32 (khối 128bit, 192 và 256bit). Tương tự với khóa, khóa cũng được biểu diễn dưới dạng byte hình chữ nhật, bao gồm 4 hàng và số cột được kí hiệu là $N_k$, bằng chiều dài của khóa chia cho 32 (khóa 128, 192 và 256bit). Số chu kỳ được chỉ định trong thuật toán sẽ phụ thuộc vào giá trị của $N_b$ và $N_k$.\quad

\begin{table}[h]
	\centering
	\begin{tabular}{|c|c|c|c|}
		\hline
		$N_r$ & $N_b = 4$ & $N_b = 6$ & $N_b = 8$ \\ \hline
		$N_k = 4$ & 10 & 12 & 14 \\ \hline
		$N_k = 6$ & 12 & 12 & 14 \\ \hline
		$N_k = 8$ & 14 & 14 & 14 \\ \hline
	\end{tabular}
\end{table}

Giải thuật AES được ứng dụng trong mã hóa \textbf{bản khối rõ}, bao gồm bốn bước xử lý chính như sau:

\begin{enumerate}[label= \textbf{Bước \arabic*:}]
	\item Mở rộng khóa nhằm thực hiện sinh các khóa chính (Round key), dùng trong các vòng lặp từ khóa chính sử dụng thủ tục sinh khóa \textbf{Rijndael}.
	\item Vòng khởi tạo sẽ thực hiện hàm AddRoundKey, trong đó mỗi byte trong State được kết hợp với khóa vòng sử dụng phép XOR.
	\item Với mỗi lặp chính, mỗi vòng sẽ bao gồm 4 hàm chính xử lý dữ liệu như sau:
	\begin{itemize}
		\item Tiến hành thay thế phi tuyến tính sử dụng hàm \textbf{SubBytes}, trong đó mỗi bytes trong state được thay thế bằng bytes khác sử dụng bảng tham chiếu S-box. Bảng tham chiếu này có thể đảo ngược và được xây dựng dựa trên hai phép biến đổi sau:
		\begin{itemize}
			\item Một phép nhân nghịch đảo trong trường hữu hạn $GF(2^8)$ được tính theo modulo một đa thức $m(x) = x^8 + x^4 + x^3 + x + 1$. $\{00\}$ được phản ánh trong chính nó.
			\item Sau đó, một phép biến đổi \textbf{Affine} (trong $GF(2)$) được áp dụng và có định nghĩa như sau:
			
			\begin{align*}
				\begin{bmatrix} 
					y_0\\y_1\\ y_2\\y_3\\y_4\\y_5\\y_6\\y_7
				\end{bmatrix} &= 
				\begin{bmatrix}
					1 & 0 & 0 & 0 & 1 & 1 & 1 & 1\\
					1 & 1 & 0 & 0 & 0 & 1 & 1 & 1\\
					1 & 1 & 1 & 0 & 0 & 0 & 1 & 1\\
					1 & 1 & 1 & 1 & 0 & 0 & 0 & 1\\
					1 & 1 & 1 & 1 & 1 & 0 & 0 & 0\\
					0 & 1 & 1 & 1 & 1 & 1 & 0 & 0\\
					0 & 0 & 1 & 1 & 1 & 1 & 1 & 0\\
					0 & 0 & 0 & 1 & 1 & 1 & 1 & 1\\
				\end{bmatrix}
				\begin{bmatrix} 
					x_0\\x_1\\x_2\\x_3\\x_4\\x_5\\x_6\\x_7
				\end{bmatrix}
				+
				\begin{bmatrix}
					1\\1\\0\\0\\0\\1\\1\\0
				\end{bmatrix}
			\end{align*}
			
		\end{itemize}
		\item Mỗi dòng trong states được dịch một số bước trong chu kỳ, sử dụng hàm \textbf{ShiftRows}.
		\item \textbf{MixColumns} làm trộn các cột trong states, kết hợp 4 bytes trong mỗi cột. 
		Các cột trạng thái được coi là đa thức trong $GF(2^8)$ và nhân theo modulo $x^4+1$ với một đa thức cố định: $c(x)='03'x^3+'01'x^2+'01'x+02$ nghịch đảo.
		
		Biến đổi này có thể được biểu diễn như một phép nhân trong ma trân, với mỗi bytes là một phần tử trong trường $GF(2^8)$. Cho $b(x) = c(x) \times a(x)$:
		\begin{center}
			$\begin{bmatrix} b_0 \\ b_1 \\b_2 \\ b_3\end{bmatrix} =
			\begin{bmatrix}
				02 & 03 & 01 & 01 \\
				01 & 02 & 03 & 01 \\
				01 & 01 & 02 & 03\\
				03 & 01 & 01 & 02 \\
			\end{bmatrix}
			\begin{bmatrix} a_0 \\ a_1 \\ a_2 \\ a_3 \\ \end{bmatrix}$
		\end{center}
		
		\item \textbf{AddRoundKeys} là hàm kết hợp state với khóa vòng, sử dụng phép XOR.
	\end{itemize}
	\item Vòng lặp cuối cùng cũng tương tự các vòng chính khác, nhưng không thực hiện bước \textbf{MixColumns}.
\end{enumerate}

